\documentclass[10pt,letterpaper,final]{report}
\usepackage[margin=.75in]{geometry}
\usepackage[utf8]{inputenc}
\usepackage{amsmath}
\usepackage{amsfonts}
\usepackage{amssymb}
\usepackage{amsthm}
\renewcommand\qedsymbol{$\blacksquare$}
\usepackage{enumerate}
\usepackage{enumitem}
\usepackage{hyperref}
\usepackage{pdfpages}
\usepackage{graphics}
\usepackage{graphicx}
\usepackage{tikz}
\usepackage{tikz-qtree}
\usepackage{forest}
\usepackage{float}
\usetikzlibrary{automata,arrows}
\usetikzlibrary{shapes.multipart, positioning}
\usepackage{mdframed}
\usepackage[
  backend=biber,
  style=numeric-comp,
  sorting=none
]{biblatex}
\addbibresource{references.bib}

\usepackage{minted}
\usemintedstyle{algol_nu}
\usepackage{clrscode3e}
\usepackage{xltabular}
\usepackage{pdflscape}
\usepackage{tcolorbox}
\usepackage{array}
\usepackage{todonotes}
\setuptodonotes{inline}  % render \todo notes as full-width blocks in the text flow, not margin boxes

\newcommand{\reason}[1]{&& \text{#1}}

% clrscode3e comments that are pushed to the right and left-aligned with
% each other. The box is flush-right and fixed-width, so its left edge --
% where the "//" begins -- sits at the same column on every line. Size
% \cmtwidth to the widest comment so the longest one just reaches the margin.
\newlength{\cmtwidth}
\settowidth{\cmtwidth}{................................................}
\newcommand{\cmt}[1]{\`\makebox[\cmtwidth][l]{\CommentSymbol\ #1}}

% Based on template from COSC-417 at Towson University, originally authored by Marius Zimand

% To input a script, adapt the following:
% \inputminted[breaklines, breakanywhere, fontsize=\footnotesize, frame=single, fontsize=\footnotesize]{python}{script.py}

\begin{document}
\begin{center}
  \fbox{
    \vbox{
      \begin{flushleft}
        Jonathan Llovet \\  % authors' names
        EN.625.252.81 \\  %class
        2026-10-15 \\  % date
        Module 7 - Homework 7 \\ % ID
      \end{flushleft}
      \center{\Large{\textbf{Linear Algebra}}}
      %\end{mdframed}
    } % end vbox
  } % end fbox
  \vline
\end{center}

\medskip
\noindent\textbf{Prompt 1:}
\medskip

Properties of determinants: Justify your answers with specific arguments or counterexamples.

\begin{enumerate}
  \item Is it true that $\det(AB) = \det(A)*\det(B)$?
  \item Is it true that $\det(A+B) = \det(A) + \det(B)$?
  \item Is it true that $\det(A) = \det(AT)$?
\end{enumerate}

\medskip
\noindent\textbf{Answer:}
\medskip

\pagebreak
\noindent\textbf{Prompt 2:}
\medskip

Define the following matrices: \emph{nilpotent}, \emph{skew-symmetric}, \emph{orthogonal}, and \emph{unitary}.  For each type of matrix, provide the definition, an example, and an application.  If possible, give a property or formula for the determinant of each matrix.

\medskip
\noindent\textbf{Answer:}
\medskip


\pagebreak
\noindent\textbf{Prompt 3:}
\medskip

Define a \emph{Group} as used in abstract algebra. Then look up what the General Linear Group $GL_n(R)$ is as well as the Special Linear Group $SL_n(R)$.  Provide a brief summary of some properties or uses of these groups.

\medskip
\noindent\textbf{Answer:}
\medskip


\pagebreak
\noindent\textbf{Prompt 4:}
\medskip

What is the complexity, or run-time, for matrix products? Does there exist an algorithm that can compute the product of two $n \times n$ matrices with complexity $O(n^2)$?

\medskip
\noindent\textbf{Answer:}
\medskip

The time complexity of naive matrix multiplication is $\Theta(n^3)$ \cite[Ch. 4.1 Multiplying square matrices, pp.80-90]{CormenEtAl2022IntroductionToAlgorithms}. While it's true that there are $n^2$ output elements, which admittedly makes it tempting to think that there might be a way to get the result in $O(n^2)$ steps, it actually takes $\Theta(n^3)$ time (meaning that it grows exactly as $n^3$). From CLRS:

\begin{tcolorbox}[colback=blue!10!white, colframe=blue!10!black, title={Naive Matrix Multiplication - \citetitle[p. 81]{CormenEtAl2022IntroductionToAlgorithms}}, top=-10pt]
  \begin{codebox}\Procname{$\proc{MATRIX-MULTIPLY}(A,B,C,n)$}
    \li \For $i = 1$ \To $n$                \cmt{compute entries in each of $n$ rows}
    \li     \Do \For $j = 1$ \To $n$        \cmt{compute $n$ entries in row $i$}
    \li         \Do \For $k = 1$ \To $n$
    \li             \Do $c_{ij} = c_{ij} + a_{ik} \cdot b_{kj}$ \cmt{add in another term of equation~(4.1)}
                    \End
                \End
            \End
  \end{codebox}
\end{tcolorbox}

Where equation (4.1) describes the value of each entry of the resulting matrix $C$ and is given as \[
c_{ij} = \sum_{k=1}^{n}a_{ik} \cdot b_{kj}
\]

It's worth noting that this naive form of the algorithm is framed as $C = C + A \cdot B$ rather than $C = A \cdot B$. This is a more general form of the problem, and although it may require some additional steps, it does not affect the asymptotic time complexity of the algorithm, since initializing the elements of $C$ takes $\Theta(n^2)$, which is dominated by $\Theta(n^3)$.

Surprisingly, there is an algorithm for matrix multiplication that has a better runtime complexity. As is discussed in CLRS \cite[Ch. 4.2, pp. 85-89]{CormenEtAl2022IntroductionToAlgorithms}, Strassen's algorithm exploits clever algebraic manipulations and a divide and conquer approach that breaks up the given matrices into submatrices to achieve a runtime of $\Theta(n^{\lg7}) = O(n^{2.81})$.

There is a lecture here \citeurl{AbdulBari2018StrassensMatrix} about the details of Strassen's algorithm, and CLRS \cite{CormenEtAl2022IntroductionToAlgorithms} contains a discussion as well.

\pagebreak

\section*{Source Code}
The full source code, including LaTeX, tooling, other artifacts, and git history is available on my Github repository for the class here:
\begin{center}
  \url{https://github.com/jllovet/jhu-en-625-252-fa26-linear-algebra}
\end{center}

\printbibliography

\end{document}
